\section{Method}
\label{sec:method}

\subsection{Problem setting}
\label{sec:problem}

A GridWorld instance is defined by an initial board $S_0 \in
\mathbb{Z}^{N\times N}$ over $C=16$ cell types (empty, walls, and agent colours)
and a sequence of $K$-agent actions $a_{1:t} \in \mathbb{Z}^{K\times t}$. The
environment is deterministic with a fixed agent priority order: at each step,
agents $1,\dots,K$ move in turn, and collisions are resolved by priority. The
learning task is to predict the full board trajectory
$S_{1:T} \in \mathbb{Z}^{T\times N\times N}$ from $(S_0, a_{1:T})$. Crucially,
the model is trained on trajectories of length $T_{\text{train}}{=}100$ and
evaluated out-of-distribution at $T_{\text{test}}{=}150$, i.e.\ $50\%$ beyond
the training horizon. All state predictions are causal in time:
$\hat{S}_t$ may depend only on $S_0$ and $a_{1:t}$.

\subsection{The unified spatiotemporal tensor}
\label{sec:tensor}

The central design choice is to keep the entire spatiotemporal state as a
single tensor $x\in\mathbb{R}^{B\times T\times N^{2}\times d}$ throughout the
network, rather than maintaining separate hidden states per chain. Each of the
three chains then scans $x$ as a different \emph{view}: the temporal chain
partitions $x$ by cell ($B N^2$ sequences of length $T$); the spatial chain
partitions $x$ by timestep ($B T$ sequences of length $N^2$); the causal chain
operates on the action embedding $a \in \mathbb{R}^{B\times K\times T\times d}$
($B T$ sequences of length $K$). Because all three chains read and write the
same tensor, cross-chain coupling is provided \emph{implicitly} by residual
fusion at every layer---an alternative to explicit base-pair rungs that we
ablate in \cref{sec:exp-ablation}.

\paragraph{Anchor initialisation.}
The tensor is initialised by summing three embeddings:
\begin{equation}
x_{b,t,i,:} \;=\; \mathrm{emb}_{\text{cell}}(S_{0,b,i})
               \;+\; \tfrac{1}{K}\textstyle\sum_{k} \mathrm{emb}_{\text{act}}(a_{b,k,t})
               \;+\; \mathrm{emb}_{\text{time}}(t),
\label{eq:anchor}
\end{equation}
where $i$ indexes the $N^2$ cells flattened in row-major order. The initial
board $S_0$ is broadcast across all $T$ timesteps (the model must learn to
\emph{modify} it forward in time), the mean action embedding injects
per-step intent, and the time embedding provides absolute position. We use
$d=256$.

\subsection{HeteroMamba2: three heterogeneous chains}
\label{sec:hetero}

The core block $\mathrm{HeteroMamba2}$ stacks $L=2$ layers, each performing
three parallel scans followed by residual fusion. The three chains have
heterogeneous Mamba2 hyperparameters chosen to match their structural
characteristics (\cref{tab:chains}).

\begin{table}[h]
\centering
\small
\caption{Per-chain Mamba2 (SSD) configuration. ``Dir.'' is the scan direction;
``Causal'' means strictly left-to-right, ``Bi'' means bidirectional (sum of
forward and flipped scans with shared parameters).}
\label{tab:chains}
\begin{tabular}{lccccc}
\toprule
Chain & Scans over & $d_{\text{state}}$ & expand & headdim & Dir. \\
\midrule
Spatial (row) & $N^2$ cells per $(b,t)$ & 128 & 1 & 32 & Bi \\
Spatial (col) & $N^2$ cells per $(b,t)$ & 128 & 1 & 32 & Bi \\
Temporal      & $T$ steps per $(b,i)$   & 64  & 2 & 64 & Causal \\
Causal        & $K$ agents per $(b,t)$  & 32  & 2 & 64 & Causal \\
\bottomrule
\end{tabular}
\end{table}

\paragraph{Spatial chain (bidirectional, row + column).}
The grid has no canonical left-to-right order over cells, so the spatial chain
is bidirectional: $y = \mathrm{Mamba2}(x) + \mathrm{flip}(\mathrm{Mamba2}(\mathrm{flip}(x)))$
with shared parameters to control the budget. Because a 2-D grid has two natural
orders (row-major and column-major), we run two independent bidirectional
Mamba2 modules---one over the row-major flattening, one over the column-major
flattening (transposed)---and fuse their outputs:
\begin{equation}
x_s \;=\; W_s\,\big[\,\mathrm{BiMamba2}_{\text{row}}(x_{\text{row}})\,;\,
                      \mathrm{BiMamba2}_{\text{col}}(x_{\text{col}})\,\big],
\label{eq:spatial}
\end{equation}
with $W_s\in\mathbb{R}^{d\times 2d}$. The spatial chain uses $d_{\text{state}}{=}128$
(the grid is large and benefits from a richer state) and \texttt{headdim=32} with
\texttt{expand=1}, which we found necessary to avoid a \texttt{causal\_conv1d}
stride-alignment error in the Mamba2 kernel.

\paragraph{Temporal chain (causal).}
For each cell $i$ independently, a causal Mamba2 scans along $t$:
$x_t = \mathrm{Mamba2}_{\text{causal}}(x_{[:,\,:,i,:]})$, reshaped to
$(B N^2, T, d)$ and back. Causality is the Mamba2 default (the convolution and
recurrence are masked), so $\hat{S}_t$ depends only on
$S_0$ and $a_{1:t}$ as required.

\paragraph{Causal chain (agent-wise).}
The action embedding $a\in\mathbb{R}^{B\times K\times T\times d}$ is scanned
causally along the $K$-agent axis, matching the environment's fixed priority
order: agent $k$'s representation is conditioned on agents $<k$. The resulting
per-agent representation is updated residually, and its mean across $K$ is
projected and broadcast back into the cell tensor as an injection:
\begin{equation}
c_{\text{inject}} \;=\; W_c\,\big(\tfrac{1}{K}\textstyle\sum_{k} \mathrm{Mamba2}_{\text{causal}}(a)\big),
\label{eq:causal}
\end{equation}
with $W_c\in\mathbb{R}^{d\times d}$, broadcast from $(B,T,1,d)$ to $(B,T,N^2,d)$.

\paragraph{Residual fusion.}
Each layer updates the unified tensor by residual addition followed by
LayerNorm:
\begin{equation}
x \;\leftarrow\; \mathrm{LN}\!\big(x + x_s + x_t + c_{\text{inject}}\big).
\label{eq:fusion}
\end{equation}
This is the single point of cross-chain interaction. We deliberately avoid
global pooling, cross-attention, or learned gates at fusion---the ablations in
\cref{sec:exp-ablation} show that adding explicit coupling here is either
harmful or neutral.

\subsection{Output head and loss}
\label{sec:head}

A simple MLP head projects each cell back to the $C=16$ cell classes:
$\mathrm{head}(x) = \mathrm{Linear}_{d\to C}(\mathrm{GELU}(\mathrm{Linear}_{d\to d/2}(\mathrm{LN}(x))))$,
yielding logits $\hat{S}\in\mathbb{R}^{B\times T\times N\times N\times C}$.

\paragraph{Balanced cross-entropy loss.}
A naive cross-entropy on all cells admits the all-zero shortcut: since the
empty cell dominates, predicting class~0 everywhere yields low loss and high
accuracy. We use a \emph{balanced} cross-entropy that up-weights cells which
actually change between $S_0$ and $S_t$, and down-weights the empty class:
\begin{equation}
\mathcal{L} \;=\; \sum_{b,t,i,j} w_{ij}^{(t)}\,w_{\text{class}}(S_{t,b,i,j})\;
                  \mathrm{CE}\!\big(\hat{S}_{t,b,i,j},\, S_{t,b,i,j}\big),
\label{eq:loss}
\end{equation}
where $w_{ij}^{(t)}{=}20$ if $S_{t,i,j}\neq S_{0,i,j}$ and $1$ otherwise, and
$w_{\text{class}}(0){=}0.1$ with $w_{\text{class}}(c){=}1$ for $c\neq 0$. The
loss is computed on the final timestep plus an auxiliary term (weight $0.3$)
on intermediate timesteps to stabilise credit assignment over the long horizon.
This loss is necessary but, as the Transformer ablation shows, not sufficient:
the same loss with the same data still collapses for an equally sized
attention model.

\subsection{Base-pair ablation variants}
\label{sec:method-bp}

To test whether explicit DNA-style base-pair rungs complement the implicit
residual coupling of \cref{eq:fusion}, we define two architectural variants
that are otherwise identical to TriHelix-Mamba2 (same tensor, same chains, same
head, same loss):

\begin{itemize}[leftmargin=*,itemsep=2pt]
  \item \textbf{BPv1 (in-chain base pair).} Each chain applies an internal
        complementary-gate plus an anti-symmetric head before its output is
        fused. This formalises base-pairing \emph{within} a strand.
  \item \textbf{BPv2 (cross-chain base pair).} A single learned rung pulls the
        three chains toward a shared state $\bar{x}$ before residual fusion:
        $x_c \leftarrow x_c + \alpha\,(\bar{x} - x_c)$ for each chain $c$, with
        $\alpha\in\mathbb{R}^{2}$ initialised to $0$ and learned per layer.
        This formalises base-pairing \emph{across} the three strands.
\end{itemize}

\noindent Both variants are evaluated in \cref{sec:exp-ablation}. The full
(non-ablated) TriHelix-Mamba2 is our main model.

\subsection{Implementation notes}
\label{sec:impl}

The model is implemented in PyTorch with the \texttt{mamba\_ssm} 2.2.4 kernel
(Mamba2/SSD). Three hyperparameter constraints from the kernel shaped the
design: (i) the 1-D convolution width \texttt{d\_conv} must be in $\{2,3,4\}$;
(ii) the spatial chain's \texttt{expand=1} requires \texttt{headdim=32} to
avoid a \texttt{causal\_conv1d} stride-alignment error
(\texttt{headdim=64} fails); (iii) \texttt{headdim} must divide
$d_{\text{model}}\times\texttt{expand}$. The total parameter count is 3.26M.
Training uses AdamW with batch size 4 on a single 8\,GB GPU
(\texttt{PYTORCH\_CUDA\_ALLOC\_CONF=expandable\_segments:True}) for 2000 steps.
